\ysrc{137}
让我们看线性情形。固定 $\mathbb Q$ 中若干斜率 $\lambda_1>\cdots>\lambda_r$，并取重数 $m_1,\cdots,m_r$。相应的向量丛是 $\mathcal E=\mathcal O(\lambda_1)^{m_1}\oplus\cdots\oplus\mathcal O(\lambda_r)^{m_r}$。记 $D_\lambda$ 为 $E$ 上具有不变量 $\lambda$ 的除代数（division algebra）。$\mathcal E$ 的自同构层是下列三者的半直积：$D^\times_{\lambda_1}\times\cdots\times D^\times_{\lambda_r}$、它的连通分支群，以及一个幂幺钻石（unipotent diamond）；该幂幺钻石是若干 Banach--Colmez 空间的逐次扩张，这些空间对应于 $\mathcal Hom(\mathcal O(\lambda_i)^{m_i},\mathcal O(\lambda_j)^{m_j})$，$i>j$。正向量丛所对应的 Banach--Colmez 空间的维数等于该向量丛的次数。于是这个幂幺钻石的维数为
\[
\sum_{i>j}m_j\deg\mathcal O(\lambda_j)-m_i\deg\mathcal O(\lambda_i).
\]

现在有下述结果。
\begin{thm}{2.20}
与 $[b]$ 相联系的 Harder--Narasimhan 分层同构于分类叠 $[*/\widehat G_b]$。\label{thm:a3-2-20}
\end{thm}

这意味着：若 $\mathcal E$ 是 $X_S$ 上的 $G$-丛，且在 $S$ 上逐几何纤维同构于 $\mathcal E_b$，则在 $S$ 上 $v$-局部地有 $\mathcal E\cong\mathcal E_b$。事实上我们甚至证明，这在 $S$ 上 pro-平展局部地已经成立。对线性群而言，这意味着：对 $X_S$ 上牛顿多边形在 $S$ 上逐纤维常值的向量丛，我们可以在 $S$ 上局部地分裂其 Harder--Narasimhan 滤链。这一结果在 [14] 中已有证明，但我们用钻石理论给出一个新证明。
\begin{cor}{2.21}
与 $[b]\in B(G)$ 相联系的 HN 分层在 $*$ 上是分离的、$\ell$-上同调光滑的，维数为 $-\langle\nu_b,2\rho\rangle$。\label{cor:a3-2-21}
\end{cor}

\subsection{光滑性的 Jacobi 判据（[9, Chapter IV.4]）}\label{sec:a3-2-7}

\subsubsection{陈述}

为在 $\Bun_G$ 上构造一些「好看」的图表，我们需要经典情形下所谓 Quot 概形的某种类比或变体。为此，取 $S\in\mathrm{Perf}_{\mathbb F_q}$，并取次完美胚（sous-完美胚）adic 空间之间的光滑态射
\[
Z\longrightarrow X_S.
\]
此外作如下「拟射影」假设：存在从 $Z$ 到 $\mathbb P^n_{X_S}$ 的某个开子集的 Zariski 闭浸入。我们可能要考虑的基本例子如下。设 $S=\Spa(R,R^+)$ 是仿射完美胚，$\mathscr X_S$ 是相应的「代数曲线」，即视作 $\Spec(E)$ 上概形的那条曲线。设 $\mathfrak Z\to\mathscr X_S$ 是概形的光滑拟射影态射。于是可以定义 $\mathfrak Z^{ad}\to\mathscr X_S$，使它对 $Z$ 满足前述假设。此外，adic 化给出双射
\[
\{\mathfrak Z\to\mathscr X_S\ \text{的截面}\}\xrightarrow{\ \sim\ }\{\mathfrak Z^{ad}\to\mathscr X_S\ \text{的截面}\}.
\]
现在定义
\[
\mathcal M_Z\longrightarrow S
\]
为「$Z\to X_S$ 的截面的模空间」，即在 $S$-完美胚空间（perfectoid space）上的函子，它把 $T/S$ 送到截面 $s$
\begin{center}
\begin{tikzcd}
& Z \arrow[d] \\
X_T \arrow[ur, "s"] \arrow[r] & X_S
\end{tikzcd}
\end{center}
\ysrc{138}
它与 $Z\times_{X_S}X_T\to X_T$ 的截面相同（该纤维积作为 $X_T$ 上的光滑次完美胚空间有意义）。可以把 $Z\to X_Z$ 的切丛 $T_{Z/X_S}$ 定义为 $Z$ 上的向量丛。\tn{底本此处作 $Z\to X_Z$，按上下文应为 $Z\to X_S$，疑为笔误。}接着定义
\[
\mathcal M_Z^{sm}\subset\mathcal M_Z
\]
为开子函子，要求经由前述截面 $s$，$s^*T_{Z/X_S}$ 是 $X_T$ 上的向量丛，且在 $T$ 上逐纤维具有正的（非零的）HN 斜率。下面是我们的光滑性 Jacobi 判据。
\begin{thm}{2.22}
(1) 函子 $\mathcal M_Z$ 可由局部空间钻石表示。

(2) 态射 $\mathcal M_Z^{sm}\to S$ 是分离的、$\ell$-上同调光滑的，其维数等于由截面 $s$ 给出的点处 $s^*T_{Z/X_S}$ 的次数。\label{thm:a3-2-22}
\end{thm}
\begin{rem}{2.23}
在「线性情形」，即 $Z$ 是 $X_S$ 上向量丛 $\mathcal E$ 的几何实现 $\mathbb V(\mathcal E)$ 时，有 $\mathcal M_Z=\mathcal{BC}(\mathcal E)$，上述结果就是 Banach--Colmez 空间理论中的基本结果。因此，前述结果把这一结果推广到了曲线上更一般的「非线性」代数方程。\label{rem:a3-2-23}
\end{rem}
\begin{exam}{2.24}
设 $n\ge1$ 是整数，$Z$ 是 Fermat 曲线 $\{[x:y:z]\in\mathbb P^2_E\mid x^n+y^n=z^n\}$ 到 $X_S$ 的拉回。有分解 $\mathcal M_Z=\coprod_{k\in\mathbb Z}\mathcal M^k_Z$，其中 $\mathcal M^k_Z$ 是使线丛 $\mathcal O_{\mathbb P^2}(1)$ 的拉回次数为 $k$ 的开/闭子叠。设 $\mathbb B$ 是环的 $v$-层 $S\mapsto\mathcal O(Y_S)$。于是当 $k<0$ 时 $\mathcal M^k_Z=\emptyset$，而当 $k\ge0$ 时
\[
\mathcal M^k_Z=\{(x,y,z)\in(\mathbb B^{\varphi=\pi^k})^3\setminus\{(0,0,0)\}\mid x^n+y^n=z^n\}/E^\times.
\]
\label{exam:a3-2-24}
\end{exam}
\begin{rem}{2.25}
在概形的「经典情形」中，前述结果众所周知且直接。更确切地说，若 $X$ 是域 $k$ 上的光滑射影曲线，$S$ 是 $k$-概形，$Z\to X_S$ 是拟射影光滑的，则按定义 $\mathcal M_Z^{sm}\to S$ 是 $Z\to X_S$ 的截面 $s$ 的模空间，要求向量丛 $s^*T_{Z/X_S}$ 在基上逐纤维没有 $H^1$。容易验证该函子是形式光滑的。\label{rem:a3-2-25}
\end{rem}

例如，我们处理 Quot 钻石的情形。设 $\mathcal E$ 是 $X_S$ 上的向量丛。令
\[
\mathrm{Quot}_{\mathcal E}\longrightarrow S
\]
为把 $T/S$ 送到 $\mathcal E|_{X_T}$ 的非零局部自由商的函子。$\mathcal E$ 的非零局部自由商的模空间 $Z\to X_S$ 是 Grassmann 簇的无交并 $\coprod_{r\ge1}\mathrm{Gr}_r(\mathcal E)\to X_S$，其中 $r$ 是商的秩。我们有 $\mathcal M_Z=\mathrm{Quot}_{\mathcal E}$。此时 $\mathcal M_Z^{sm}\to S$ 把 $T/S$ 送到局部自由商 $u:\mathcal E|_{X_T}\twoheadrightarrow\mathcal F$，使得 $(\ker u)^\vee\otimes\mathcal F$（按前述记号即 $s^*T_{Z/X_S}$）在 $T$ 上逐纤维具有正的 HN 斜率，也就是说 $\ker u$ 的最大斜率严格小于 $\mathcal F$ 的最小斜率。按 Jacobi 判据，这是分离的、$\ell$-上同调光滑的局部空间钻石。

\subsubsection{证明的若干工具：ULA 层}

\ysrc{139}
设 $\Lambda$ 是挠 $\mathbb Z_\ell$-代数。设 $f:X\to Y$ 是局部空间钻石之间的可紧化态射，且局部具有有限 dim.trg.\ 维数。设 $A\in D_{\text{ét}}(X,\Lambda)$（关于 $D_{\text{ét}}(X,\Lambda)$ 的定义，见 §3.1 开头的讨论）。我们定义 $A$ 为 $f$-ULA（universally locally acyclic，泛局部无环）这一概念，它满足下列性质（[9, Chapter IV.2]）。这些性质与概形情形中的相应性质类似：\tn{原文缩略语 “dim.trg.” 未展开，照原样保留；上下文推断指该态射局部上具有有限维的相对结构，无外部依据。}
\begin{enumerate}[label=(\arabic*),leftmargin=2.4em,itemsep=0.2em]
\item 若 $f$ 是 $\ell$-上同调光滑的，且 $A$ 局部常值、具有完美纤维，则 $A$ 是 $f$-ULA 的。
\item 若 $f$ 是恒等态射，则 $A$ 是 $f$-ULA 的当且仅当 $A$ 局部常值、具有完美纤维。
\item 若有图表 $X'\xrightarrow{g}X\xrightarrow{f}Y$，其中 $g$ 分离、$\ell$-上同调光滑且满，则 $A$ 是 $f$-ULA 的当且仅当 $g^*A$ 是 $f\circ g$-ULA 的：ULA 这一概念在「源上光滑局部」。
\item 若有图表 $X'\xrightarrow{g}X\xrightarrow{f}Y$，其中 $g$ 固有，且 $A=Rg_*B$，$B$ 是 $f\circ g$-ULA 的，则 $A$ 是 $f$-ULA 的。
\end{enumerate}

尽管这不是我们采用的定义，ULA 层最快的形式化定义是借助 Lu--Zheng [19] 关于可对偶性的工作的类比（见 [9, Chapter IV.2.3.3]）：即上同调对应的 2-范畴中的左伴随。更确切地说，若 $Y$ 是固定的局部空间钻石，则对应的 2-范畴 $\mathcal C_Y$ 如下：
\begin{itemize}[leftmargin=1.6em,itemsep=0.2em]
\item 其对象是可紧化、局部有限 dim.trg.\ 的局部空间钻石的态射 $X\to Y$；
\item $X\to Y$ 与 $X'\to Y$ 之间的 1-态射由 $D_{\text{ét}}(X\times_Y X',\Lambda)$ 的对象给出，我们把它们视作核；对 $A\in D_{\text{ét}}(X\times_Y X',\Lambda)$，相应的变换是
\[
F_A:\ D_{\text{ét}}(X,\Lambda)\longrightarrow D_{\text{ét}}(X',\Lambda),\qquad
\mathcal F\longmapsto pr_{2!}(pr_1^*\mathcal F\otimes^L_\Lambda A);
\]
\item 1-态射的复合由「核的卷积」给出，因此 $F_{A_1*A_2}=F_{A_1}\circ F_{A_2}$；
\item $D_{\text{ét}}(X\times_Y X',\Lambda)$ 中 $A_1$ 与 $A_2$ 之间的 2-态射就是 $A_1$ 与 $A_2$ 之间的态射。
\end{itemize}

现在可以回到 $f$-ULA 平展复形的定义。任何 $A\in D_{\text{ét}}(X,\Lambda)$ 都看作从 $X\to Y$ 到 $Y\xrightarrow{\mathrm{Id}}Y$ 的 1-态射。若该 1-态射在 2-范畴 $\mathcal C_Y$ 中是左伴随，则称 $A$ 为 $f$-ULA 的（任意 2-范畴中左伴随的定义，是范畴的 2-范畴中通常定义的自然推广）。于是可以证明：这个左伴随只要存在，就自动是其 Verdier 对偶 $\mathbb D_{X/Y}(A)$。

ULA 层在 Verdier 对偶与基变换下的性质很好。我们证明：
\begin{enumerate}[label=(\arabic*),leftmargin=2.4em,itemsep=0.2em]
\item 若 $A$ 是 $f$-ULA 的，则 $\mathbb D_{X/Y}(A)$ 也是 $f$-ULA 的，此时
\[
A\xrightarrow{\ \sim\ }\mathbb D_{X/Y}(\mathbb D_{X/Y}(A)).
\]
\item 若 $g:Y'\to Y$ 是局部空间钻石的态射，$\tilde g:X\times_Y Y'\to X$，则当 $A$ 是 $f$-ULA 时，
\[
\tilde g^*\mathbb D_{X/Y}(A)\xrightarrow{\ \sim\ }\mathbb D_{X\times_Y Y'/Y'}(\tilde g^*A).
\]
\ysrc{140}
\item 若 $A$ 是 $f$-ULA 的，则对任意 $B\in D_{\text{ét}}(Y,\Lambda)$ 有
\[
\mathbb D_{X/Y}(A)\otimes^L_\Lambda f^*B\xrightarrow{\ \sim\ }R\mathcal Hom_\Lambda(A,Rf^!B).
\]
\end{enumerate}
\begin{rem}{2.26}
我们给出的 $f$-ULA 平展复形定义实际上受性质 (3) 约束。事实上，对任意 $B\in D_{\text{ét}}(Y,\Lambda)$，性质 (3) 蕴含
\[
R\Hom(Rf_!A,B)=R\Gamma(X,\mathbb D_{X/Y}(A)\otimes^L_\Lambda f^*B).
\]
此外，若 $f$ 是拟紧的，则该性质蕴含：$Rf_!A$ 限制到 $Y$ 的任意拟紧开子集 $U$ 上，是 $D_{\text{ét}}(U,\Lambda)$ 的紧对象。这等价于说（[23, Chapter 20]）$Rf_!A$ 可构造，且具有完美的几何茎。我们采用的定义正是如此：在 $X$ 上平展局部地强制这一可构造性（因为 ULA 概念必须在源上平展局部）。\label{rem:a3-2-26}
\end{rem}

对 $A=\mathbb F_\ell$ 应用性质 (2) 与 (3)，得到下面结果。
\begin{thm}{2.27}
态射 $f$ 是 $\ell$-上同调光滑的，当且仅当

(1) $\mathbb F_\ell$ 是 $f$-ULA 的；

(2) $Rf^!\mathbb F_\ell$ 可逆，即局部同构于 $\mathbb F_\ell[d]$，其中 $d\in\mathbb Z$。\label{thm:a3-2-27}
\end{thm}

在光滑性 Jacobi 判据的证明中，我们正是把证明切成两部分来使用这一点。

\subsubsection{证明的若干工具：形式光滑性与形变到法锥}

定理 \ref{thm:a3-2-22} 的证明分两部分，第一部分在于证明 $\mathbb F_\ell$ 关于 $\mathcal M_Z^{sm}\to S$ 是 ULA 的。这用到一个事实：与上同调光滑性不同，ULA 性在收缩（retract）下稳定。用一个例子说明。设 $S$ 是仿射完美胚。设 $\mathbb B^d_S\to S$ 是 $S$ 上的闭 $d$ 维完美胚球。设 $Z\subset\mathbb B^d_S$ 是 Zariski 闭的。假设存在 $Z$ 的平展邻域 $U\to\mathbb B^d_S$，以及包含映射 $Z\times_{\mathbb B^d_S}U\hookrightarrow U$ 的一个收缩。此时 $\mathbb F_\ell$ 关于 $Z\to S$ 是 $f$-ULA 的。事实上我们证明下述定理。
\begin{thm}{2.28}
态射 $\mathcal M_Z^{sm}\to S$ 是形式光滑的，意思是：对任意图表
\begin{center}
\begin{tikzcd}
T' \arrow[r] \arrow[d] & \mathcal M_Z^{sm} \arrow[d] \\
T \arrow[r] \arrow[ur, dashed, "s"] & S
\end{tikzcd}
\end{center}
其中 $T$ 是仿射完美胚、$T'$ 在 $T$ 中 Zariski 闭，把 $T$ 换成 $T'$ 的平展邻域 $U\to T$、把 $T'$ 换成 $T'\times_T U$ 之后，存在补全该图表的态射 $s$。\label{thm:a3-2-28}
\end{thm}

一个初等的论证随后表明，这蕴含 $\mathbb F_\ell$ 关于 $\mathcal M_Z^{sm}\to S$ 是 ULA 的。现在用到一个事实：若 $\mathbb F_\ell$ 是 $f$-ULA 的，则 $Rf^!\mathbb F_\ell$ 的构造与基变换交换。利用这一点，经过换基 $S$，问题归约为证明：若
\begin{center}
\begin{tikzcd}
\mathcal M_Z^{sm} \arrow[d, "f"'] \\
S \arrow[u, "i"']
\end{tikzcd}
\end{center}
\ysrc{141}
则 $i^*Rf^!\mathbb F_\ell$ 可逆。这样的截面 $i$ 对应于一个截面
\begin{center}
\begin{tikzcd}
Z \arrow[d] \\
X_S \arrow[u, "k"']
\end{tikzcd}
\end{center}
我们于是对正则浸入 $X_S\hookrightarrow Z$ 使用形变到法锥：把 $Z$ 换为该法锥形变 $C$，得到图表
\begin{center}
\begin{tikzcd}
\mathcal M_C \arrow[d] \\
S\times E \arrow[u]
\end{tikzcd}
\end{center}
它在 $0\in E$ 处的纤维是 $\mathcal{BC}(k^*T_{Z/X_S})\to S$ 的零截面，而在 $0\in E$ 之外同构于 $\mathcal M_Z\times E^\times$。该图表是 $E^\times$-等变的；利用这个作用以及 $\mathcal{BC}(k^*T_{Z/X_S})\to S$ 的上同调光滑性，即可得结论。这里 $\mathcal{BC}$ 表示 Banach--Colmez 空间。

\subsection{$\Bun_G$ 上一些「好看」的图表（[9, Chapter V.3]）}\label{sec:a3-2-8}

与 Beauville--Laszlo 单值化、或前述 Quot 钻石（它们只在线性群情形下好用）不同，我们在研究 $\Bun_G$ 上的层时使用另一些图表。它们由下述方式给出。为简化，设 $G$ 拟分裂。设 $[b]\in B(G)$，$M$ 是以 $[\nu_b]\in X_*(A)^+_{\mathbb Q}$ 的中心化子为标准 Levi 子群的那个标准 Levi 子群。模去 $\sigma$-共轭，可以假设 $b$ 是 $M(\breve E)$ 中的基本元素 $b_M$。设 $P$ 是与 $[\nu_b]$ 相配、以 $M$ 为标准 Levi 子群的标准抛物子群。定义
\[
\mathcal M_b
\]
为这样的模空间：它由 $X_S$ 上的 $P$-丛 $\mathcal E_P$ 给出，使得 $\mathcal E_P\times^P M$（Levi 商）在 $S$ 上逐几何纤维同构于 $\mathcal E_{b_M}$。于是有图表
\begin{center}
\begin{tikzcd}
\widehat{\mathcal M}_b \arrow[r] \arrow[d] & \mathcal M_b \arrow[r, "\pi_b"] \arrow[d, "q_b"'] & \Bun_G \\
* \arrow[r] & {[*/\widehat G_b(E)]} \arrow[u, "s_b"] &
\end{tikzcd}
\end{center}
其中 $q_b$ 把 $\mathcal E_P$ 送到 $\mathcal E_P\times^P M$，$\pi_b$ 送到 $\mathcal E_P\times^P G$，截面 $s_b$ 把 $\mathcal E_M$ 送到 $\mathcal E_M\times^M P$。下面这个结果是我们研究 $\Bun_G$ 上的层的核心。
\begin{thm}{2.29}
(1) 态射 $q_b:\mathcal M_b\to[*/\widehat G_b(E)]$ 可由局部空间钻石表示，且是分离的、$\ell$-上同调光滑的，维数为 $\langle\nu_b,2\rho\rangle$。特别地，$\mathcal M_b$ 是上同调光滑的 Artin $v$-叠，维数为 $\langle\nu_b,2\rho\rangle$。

(2) 态射 $\pi_b:\mathcal M_b\to\Bun_G$ 部分固有（partially proper），可由局部空间钻石表示，分离、$\ell$-上同调光滑，其像集是 $|\Bun_G|$ 中 $[b]$ 的所有泛化点构成的集合。

(3) $\widehat{\mathcal M}_b\setminus\{*\}$ 是空间钻石。\label{thm:a3-2-29}
\end{thm}
\ysrc{142}
第 (1) 点由下述事实推出：$q_b$ 是负 Banach--Colmez 空间中的迭代纤维化，见 §2.3。第 (2) 点是光滑性 Jacobi 判据（定理 \ref{thm:a3-2-22}）的直接应用。给定 $X_S$ 上的 $G$-丛 $\mathcal E$，把它视作平展 $G$-挠子 $\mathcal E\to X_S$，我们对 $Z=P\backslash\mathcal E$ 应用 Jacobi 判据。此时空间 $\mathcal M_Z\to S$ 就是 $\mathcal E$ 到 $P$ 的约化的模空间，即 $\Bun_P\times_{\Bun_G}S$。第 (3) 点类似于「穿孔的绝对负 Banach--Colmez 空间是空间钻石」的证明，见 §2.3.2。
\begin{rem}{2.30}
$\widehat{\mathcal M}_b\setminus\{*\}$ 的空间性、以及它是拟紧的（而不只是局部空间的）这一事实，是我们工作中的关键工具。这也是我们宁愿「绝对地」在 $*$ 上使用 $\Bun_G\to*$、而不用标量扩张到某个代数闭完美胚域 $C$ 之后的 $\Bun_G\times\Spa(C)$ 的主要原因之一。事实上，$\widehat{\mathcal M}_b\times\Spa(C)$ 是局部空间钻石，但 $(\widehat{\mathcal M}_b\setminus\{*\})\times\Spa(C)$ 不再是拟紧的。这类现象在我们的世界中经常出现，其根源在于 $v$-意象的终对象 $*$ 不是拟分离的。例如，带截面 $\{*\}\hookrightarrow D$（由 $x_1=\cdots=x_n=0$ 给出）的 $v$-层
\[
D=\Spa\big(\overline{\mathbb F}_q[\![x_1^{1/p^\infty},\dots,x_n^{1/p^\infty}]\!],\overline{\mathbb F}_q[\![x_1^{1/p^\infty},\dots,x_n^{1/p^\infty}]\!]\big)
\]
满足：(1) $D$ 不是钻石；(2) $D\setminus\{*\}$ 是空间钻石；(3) 标量扩张到 $C$ 之后，$D\times\Spa(C)$ 是局部空间钻石，但 $(D\setminus\{*\})\times\Spa(C)$ 不是空间的。在我们的情形中情况更糟：$v$-层 $\widehat{\mathcal M}_b$ 一般甚至不能由形式概形表示。\label{rem:a3-2-30}
\end{rem}

我们处理 $G=\mathrm{GL}_2$、$b=\begin{pmatrix}1&0\\0&\pi^{-1}\end{pmatrix}$ 的情形。层 $\widehat{\mathcal M}_b$ 是 $\mathcal O(1)$ 被 $\mathcal O$ 的扩张的模空间，即
\[
\widehat{\mathcal M}_b=\mathcal{BC}(\mathcal O(-1)[1])\longrightarrow *.
\]
我们有
\[
\mathcal M_b=[\mathcal{BC}(\mathcal O(-1)[1])/E^\times\times E^\times],
\]
其中 $(a,b)\in E^\times\times E^\times$ 通过 $ab^{-1}$ 作用在 Banach--Colmez 空间上。态射
\[
[\mathcal{BC}(\mathcal O(-1)[1])/E^\times\times E^\times]\longrightarrow\Bun_{\mathrm{GL}_2}
\]
把截面 $[*/E^\times\times E^\times]$ 送到 $|\Bun_G|$ 中由 $[b]$ 给出的点。其补集送到与 $\begin{pmatrix}0&\pi^{-1}\\1&0\end{pmatrix}$ 相联系的半稳定点，即 $\mathcal O(1)$ 被 $\mathcal O$ 的扩张逐几何纤维同构于 $\mathcal O(\frac12)$ 的轨迹。

\section{范畴 $D_{lis}(\Bun_G,\Lambda)$}

\subsection{挠情形}

设 $\Lambda$ 是 $\mathbb Z_\ell$-代数。本节假设 $\Lambda$ 是挠的。定义
\[
D_{lis}(\Bun_G,\Lambda):=D_{\text{ét}}(\Bun_G,\Lambda).
\]
由于尚未解释，下面说明挠 $\Lambda$ 时 $D_{\text{ét}}(\Bun_G,\Lambda)$ 的定义。该定义在 [23] 中分几步、通过对严格全不连通情形的下降程序给出；在那种情形下，定义没有选择余地：
\ysrc{143}
\begin{enumerate}[label=(\arabic*),leftmargin=2.4em,itemsep=0.2em]
\item 完美胚空间的平展态射有好的概念。对完美胚空间 $X$，有平展 site $X_{\text{ét}}$ 及相应的意象。此时 $D^+_{\text{ét}}(X,\Lambda)$ 就是 $X_{\text{ét}}$ 上 $\Lambda$-模层的通常导出范畴。
\item 完美胚空间的分离平展态射沿 pro-平展覆盖（甚至沿 $v$-覆盖）下降。于是可以定义钻石的（局部分离）态射这一概念。对钻石 $X$，由此可以定义平展 site $X_{\text{ét}}$，并令 $D^+_{\text{ét}}(X,\Lambda)=D^+(X_{\text{ét}},\Lambda)$。
\item 对钻石 $X$，用 $D^+_{\text{ét}}(X,\Lambda)$ 记以 $D^+_{\text{ét}}(X,\Lambda)$ 为同伦范畴的那个自然稳定 $\infty$-范畴。可以证明 $X\mapsto D^+_{\text{ét}}(X,\Lambda)$ 是 $v$-超层。作为推论，对钻石 $X$ 得到 $D^+_{\text{ét}}(X,\Lambda)$ 的如下描述：
\[
\begin{aligned}
D^+_{\text{ét}}(X,\Lambda)=\Big\{&A\in\underbrace{D^+(X_v,\Lambda)}_{v\text{-层，取值于 }\Lambda\text{-模}}\Big|\ \forall f:S\to X,\ S\ \text{严格全不连通},\\
&f^*A\in D^+_{\text{ét}}(S,\Lambda)=D^+(|S|,\Lambda)\Big\}.
\end{aligned}
\]
\item 一般地，对小 $v$-叠 $X$ 定义
\[
D_{\text{ét}}(X,\Lambda)=\varprojlim_{S\to X\ \text{严格全不连通}}\underbrace{D_{\text{ét}}(S,\Lambda)}_{D(|S|,\Lambda)},
\]
其中 $S$ 取遍严格全不连通的完美胚空间；其同伦范畴不是别的，正是
\[
\begin{aligned}
D_{\text{ét}}(X,\Lambda)=\Big\{&A\in\underbrace{D(X_v,\Lambda)}_{v\text{-层，取值于 }\Lambda\text{-模}}\Big|\ \forall f:S\to X,\ S\ \text{严格全不连通},\\
&f^*A\in D_{\text{ét}}(S,\Lambda)=D(|S|,\Lambda)\Big\}.
\end{aligned}
\]
\item 可以验证：对局部空间钻石 $X$，前面 $D_{\text{ét}}(X,\Lambda)$ 的定义给出 $D(X_{\text{ét}},\Lambda)$ 的左完备化。在有限上同调假设下——典型地，$X$ 的任意拟紧开子集的 $\ell$-上同调维数有限——有 $D_{\text{ét}}(X,\Lambda)=D(X_{\text{ét}},\Lambda)$，即 $D(X_{\text{ét}},\Lambda)$ 左完备；但一般情形下这未必成立。
\end{enumerate}

于是具体地说，若 $S_\bullet\to\Bun_G$ 是由严格全不连通完美胚空间的无交并给出的 $v$-超覆盖，则
\[
D_{\text{ét}}(\Bun_G,\Lambda)=\underbrace{D(|S_\bullet|,\Lambda)}_{\text{笛卡尔层的导出范畴}}.
\]
\begin{rem}{3.1}
我们之所以要在左右两侧都处理无界的平展复形，理由如下。$v$-叠 $\Bun_G$ 不是拟紧的，因此必须在 $D^+_{\text{ét}}(\Bun_G,\Lambda)$ 中与可能无界的复形打交道。但我们又想研究这些复形的 Verdier 对偶，这就引出 $D_{\text{ét}}(\Bun_G,\Lambda)$ 的引入与研究。\label{rem:a3-3-1}
\end{rem}

分析 $D_{\text{ét}}(\Bun_G,\Lambda)$ 之对象的一种方式如下。对 $[b]\in B(G)$，有相应 HN 分层的包含
\[
i^b:[*/\widehat G_b]\hookrightarrow\Bun_G.
\]
我们想通过 $i^b$ 拉回 $\Bun_G$ 上的平展层，以理解这些层。为此需要下述结果，见 [9, Chapter V.2]。
\ysrc{144}
\begin{thm}{3.2}
对 $[b]\in B(G)$，有等同
\[
D(G_b(E),\Lambda)=D_{\text{ét}}([*/\widehat G_b(E)],\Lambda)=D_{\text{ét}}([*/\widehat G_b],\Lambda),
\]
其中左边的范畴是 $G_b(E)$ 的、以 $\Lambda$ 为系数的光滑表示的导出范畴。\label{thm:a3-3-2}
\end{thm}

第一个等同是容易的。第二个更微妙，本质地用到 $\ell\ne p$ 这一事实。更确切地说，回忆 $\widehat G_b=\widehat G_b^0\rtimes G_b(E)$。此外 $\widehat G_b^0$ 是正 Banach--Colmez 空间的逐次扩张，因而是「$\ell$-平展可缩的」（我们给出这一说法的精确含义）。第二个等同正来源于此。

最后，对每个 $[b]\in B(G)$ 有函子
\begin{center}
\begin{tikzcd}
D_{\text{ét}}(\Bun_G,\Lambda) \arrow[r, shift left=1.2pt, "(i^b)^*"] & D(G_b(E),\Lambda) \arrow[l, shift left=1.2pt, "(i^b)_!"]
\end{tikzcd}
\end{center}
它给出三角范畴 $D_{\text{ét}}(\Bun_G,\Lambda)$ 关于族 $\{D(G_b(E),\Lambda)\}_{[b]\in B(G)}$ 的半正交分解。特别地，经由 $(i^1)^*$ 与 $(i^1)_!$，范畴 $D(G(E),\Lambda)$ 是 $D_{\text{ét}}(\Bun_G,\Lambda)$ 的直因子。这意味着「$G(E)$ 的、以 $\Lambda$ 为系数的经典光滑表示论是 $D_{\text{ét}}(\Bun_G,\Lambda)$ 的直因子」。这正是我们构造 L-参数时要用的。

\subsection{一般情形（[9, Chapter VII]）}

对任意 $\mathbb Z_\ell$-代数 $\Lambda$，情形更复杂。我们只解释问题并勾画解法。首先，我们不想假设 $\Lambda$ 是 $\ell$-进完备的，因为那会迫使我们在最终构造一个指向 $G$ 的 Bernstein 中心的 $\ell$-进完备化的态射，而这并非我们想做的事。即使假设 $\Lambda$ 是 $\ell$-进完备的，若考虑 $\ell$-进完备的层，即当 $n$ 变动时 $\Lambda/\ell^n\Lambda$ 上挠平展层的极限，并形式地反演 $\ell$，最终仍会落到表示论一侧：$G(E)$ 的、取值于 $\mathbb Q_\ell$-Banach 空间且具有不变格的连续表示。这不是我们想要的。我们要处理 $G(E)$ 的纯粹「代数」光滑表示。

解法来自 pro-平展固体层（pro-étale solid sheaf）的理论。设 $\Lambda$ 是任意 $\mathbb Z_\ell$-代数。对任意 Artin $v$-叠 $X$，定义
\[
D_{\text{proét}}(X,\Lambda^{\blacksquare}).
\]
我们在这一框架下发展 $\Lambda$-模的固体 pro-平展层理论，特别是五函子形式 $(f^*,Rf_*,f_\natural,R\Hom,\overset{\blacksquare}{\otimes}{}^L_\Lambda)$，其中 $f_\natural$ 是相对同调函子，定义为 $f^*$ 的左伴随（此框架中没有好的 $Rf_!$ 概念，我们用这个相对同调作为替代）。现在的主要问题如下。有函子
\[
(i^1)^*:\ D_{\text{proét}}(\Bun_G,\Lambda^{\blacksquare})\longrightarrow D(G(E),\Lambda^{\blacksquare}),
\]
其中右边的范畴是 $G(E)$（作为凝聚群）在固体 $\Lambda^\blacksquare$-模中的表示的导出范畴。这里 $\Lambda^\blacksquare$ 是凝聚环 $\Lambda^{disc}\otimes_{\mathbb Z_\ell^{disc}}\mathbb Z_\ell$。有包含
\ysrc{145}
\[
D(G(E),\Lambda)\subset D(G(E),\Lambda^{\blacksquare})
\]
它把光滑表示 $\pi$ 送到 $\pi^{disc}\otimes_{\mathbb Z_\ell^{disc}}\mathbb Z_\ell$。但范畴 $D(G(E),\Lambda^{\blacksquare})$ 比 $D(G(E),\Lambda)$ 大得多。典型地，当 $\Lambda=\mathbb Q_\ell$ 时，$G(E)$ 的任意 $\mathbb Q_\ell$-Fréchet 连续表示都给出 $G(E)$（作为凝聚群）的一个 $\mathbb Q_\ell^{\blacksquare}$-表示。

该问题的解法是定义
\[
D_{lis}(X,\Lambda)\subset D_{\text{proét}}(X,\Lambda^{\blacksquare})
\]
为包含所有 $f_\natural\Lambda$ 的、对直和稳定的最小三角子范畴，其中 $f:Y\to X$ 取遍分离的、可由局部空间钻石表示且 $\ell$-上同调光滑的态射。我们证明这给出一个「好的」三角范畴 $D_{lis}(\Bun_G,\Lambda)$，它容许由 $D(G_b(E),\Lambda)$（$[b]\in B(G)$）给出的半正交分解。特别地，$D(G(E),\Lambda)$ 是 $D_{lis}(\Bun_G,\Lambda)$ 的直因子。

\subsection{紧生成元（[9, Chapter V.4]）}

为简化陈述，本节假设 $\Lambda$ 是挠的。众所周知，$D(G(E),\Lambda)$ 是紧生成的。一组生成元由光滑表示
\[
c\text{-Ind}_K^{G(E)}\Lambda
\]
给出，其中 $K$ 是 $G(E)$ 的紧开 pro-$p$ 子群。从几何观点看，对这样的 $K$，有 Artin $v$-叠的态射
\[
f_K:[*/K]\longrightarrow[*/G(E)].
\]
于是
\[
(f_K)_!\Lambda\in D_{lis}([*/G],\Lambda)=D(G(E),\Lambda)
\]
对应于 $c\text{-Ind}_K^{G(E)}\Lambda$。

借助 §\ref{sec:a3-2-8} 的图表 $\pi_b:\mathcal M_b\to\Bun_G$，这个构造可以推广。回忆 $\mathcal M_b=[\widehat{\mathcal M}_b/G_b(E)]$。对 $G_b(E)$ 中的紧开 pro-$p$ 子群 $K$，定义
\[
f_{b,K}:[\widehat{\mathcal M}_b/K]\longrightarrow\Bun_G.
\]
由定理 \ref{thm:a3-2-29}，它是分离的、$\ell$-上同调光滑的。现在定义
\[
A^b_K=Rf_{b,K!}Rf^!_{b,K}\Lambda\in D_{lis}(\Bun_G,\Lambda).
\]
利用 $\widehat{\mathcal M}_b\setminus\{*\}$ 是空间钻石这一事实，可以证明下述定理。
\begin{thm}{3.3}
对 $A\in D_{\text{ét}}(\widehat{\mathcal M}_b,\Lambda)$，若 $i:*\hookrightarrow\widehat{\mathcal M}_b$ 是闭点，则
\[
R\Gamma(\widehat{\mathcal M}_b,A)\xrightarrow{\ \sim\ }i^*A\in D(\Lambda).
\]
\label{thm:a3-3-3}
\end{thm}

由此可以证明下述关键结果。
\begin{thm}{3.4}
对象族 $(A^b_K)_{[b]\in B(G),\,K\subset G_b(E)}$ 是 $D_{lis}(\Bun_G,\Lambda)$ 的一组紧生成元。\label{thm:a3-3-4}
\end{thm}

事实上，对 $B\in D_{lis}(\Bun_G,\Lambda)$，由于 $f^b_K$ 是 $\ell$-上同调光滑的，
\begin{align*}
R\Hom(A^b_K,B)&=R\Hom(Rf^!_{b,K}\Lambda,Rf^!_{b,K}B)\\
&=R\Gamma(\widehat{\mathcal M}_b,Rf^*_{b,K}B)\\
&=[(i^b)^*B]^K\in D(\Lambda);
\end{align*}
利用定理 \ref{thm:a3-3-3}，定理 \ref{thm:a3-3-4} 便可由该公式容易地推出。\tn{底本此处作 $Rf^{*}_{b,K}B$；按 $f_!\dashv f^!$ 的伴随性，应为 $Rf_{b,K*}B$，疑为下标与上标互倒之误。}
\ysrc{146}
这些紧生成元是我们用于下述内容的主要工具。

\subsection{有限型、容许、Bernstein--Zelevinsky 对合（[9, Chapter V.4, V.5, V.7]）}

这项工作的一条座右铭是：最终，局部 Langlands 对应中出现的自然对象不是 $G(E)$ 的光滑表示，而是 $D_{lis}(\Bun_G,\Lambda)$ 的对象。下述结果支持这一点：它说明有限型、容许这两个概念以及 Bernstein--Zelevinsky 对合都可延拓到 $D_{lis}(\Bun_G,\Lambda)$。
\begin{thm}{3.5}
设 $A\in D_{lis}(\Bun_G,\Lambda)$。

(1) $A$ 紧，当且仅当它具有有限支撑，且对一切 $[b]\in B(G)$，$(i^b)^*A\in D(G_b(E),\Lambda)$ 紧。

(2) $A$ 关于态射 $\Bun_G\to*$ 是 ULA 的，当且仅当对一切 $[b]\in B(G)$ 与一切紧开 pro-$p$ 子群 $K\subset G_b(E)$，$[(i^b)^*A]^K\in D(\Lambda)$ 是完美复形。

(3) 存在对合
\[
\mathbb D_{BZ}:\big(D_{lis}(\Bun_G,\Lambda)^\omega\big)^{op}\xrightarrow{\ \sim\ }D_{lis}(\Bun_G,\Lambda)^\omega,
\]
它延拓 $D(G(E),\Lambda)$ 上通常的 Bernstein--Zelevinsky 对合。\label{thm:a3-3-5}
\end{thm}

这里 $D(G_b(E),\Lambda)$ 中的紧性，等价于属于由 $c\text{-Ind}_K^{G_b(E)}\Lambda$（$K$ 取遍 $G_b(E)$ 的 pro-$p$ 紧开子群）生成的厚三角子范畴。
\begin{rem}{3.6}
设 $\Lambda$ 是 $\overline{\mathbb Q}_\ell$ 或 $\mathbb Q_\ell$。

(1) $G_b(E)$ 的光滑表示范畴具有有限上同调维数（这归功于 Bernstein，例如见 [22]）。因此此时紧对象就是 $D^b_{ft}(G_b(E),\Lambda)$ 的对象（有限型上同调）。于是 $A$ 紧，当且仅当它具有有限支撑，且对一切 $[b]$，$(i^b)^*A$ 是上同调为有限型的有界复形。

(2) $A$ 是 ULA 的，当且仅当对一切 $[b]$，$(i^b)^*A$ 是具有容许上同调的复形，且对一切紧开 pro-$p$ 的 $K$，$[(i^b)^*A]^K$ 有界。\label{rem:a3-3-6}
\end{rem}

前述定理的关键工具，是上一节显式给出的紧生成元 $A^b_K$。
\begin{exam}{3.7}
设 $\pi$ 是 $G(E)$ 的、以 $\overline{\mathbb Q}_\ell$ 为系数的光滑容许表示。设 $\mathcal F_\pi$ 是 $[*/G(E)]$ 上相应的层。则 $(i^1)_!\mathcal F_\pi$ 是 ULA 的，因此其 Verdier 对偶也是。由此推出 $R(i^1)_*\mathcal F_\pi$ 在所有 $[b]$ 处的茎都是具有容许上同调的复形，这是一个非平凡的有限性陈述。\label{exam:a3-3-7}
\end{exam}

\section{几何 Satake 对应（[9, Chapter VII]）}

\subsection{局部 Hecke 叠}

固定整数 $d\ge1$。令
\[
\mathrm{Div}^d\longrightarrow *
\]
为曲线上 $d$ 次有效 Cartier 除子的层。更确切地说，$\mathrm{Div}^d(S)$ 是二元组 $(\mathcal L,u)$ 的等价类集合，其中 $\mathcal L$ 是 $X_S$ 上在 $S$ 上逐纤维为 $d$ 次的线丛，$u\in H^0(X_S,\mathcal L)$ 在 $S$ 上逐纤维非零。有
\ysrc{147}
\[
\mathrm{Div}^1=\mathrm{Spd}(\breve E)/\varphi^{\mathbb Z}.
\]
该等同由态射 $\mathrm{Spd}(\breve E)\to\mathrm{Div}^1$ 推出：它把 $S$ 的一个直到化（until）$S^\sharp$ 送到相应的 Cartier 除子 $S^\sharp\hookrightarrow X_S$。此外有同构
\[
(\mathrm{Div}^1)^d/\mathfrak S_d\xrightarrow{\ \sim\ }\mathrm{Div}^d,
\]
其中该商是 pro-平展商，态射由把 $d$ 个 1 次除子相加为 1 个 $d$ 次除子给出。看 $\mathrm{Div}^d$ 的另一种方式，是把它看作穿孔的绝对 Banach--Colmez 空间的商（见 §2.3.2 与 [7]）
\[
\mathcal{BC}(\mathcal O(d))\setminus\{0\}/E^\times.
\]

存在 Beilinson--Drinfeld 型的仿射 Grassmann 簇，即 $v$-层
\[
\mathrm{Gr}_{G,\mathrm{Div}^d}\longrightarrow\mathrm{Div}^d,
\]
它在 $S$ 上的取值由以下数据给出：$X_S$ 上的 $d$ 次除子 $D$、$X_S$ 上的 $G$-丛 $\mathcal E$，以及平凡 $G$-丛与 $\mathcal E$ 在 $X_S\setminus D$ 上沿 $D$ 亚纯的同构。§2.4 与 [25] 中「通常的」$B_{dR}$ 仿射 Grassmann 簇是
\[
\mathrm{Gr}_{G,\mathrm{Div}^1}\times_{\mathrm{Div}^1}\mathrm{Spd}(\breve E).
\]
Grassmann 簇 $\mathrm{Gr}_{G,\mathrm{Div}^d}\to\mathrm{Div}^d$ 带有
\[
L^+_{\mathrm{Div}^d}G\longrightarrow\mathrm{Div}^d
\]
的作用，即相应的正圈群。按定义，对仿射完美胚 $S$，$L^+_{\mathrm{Div}^d}G(S)$ 由 $D\in\mathrm{Div}^d(S)$ 与
\[
\varprojlim_{k\ge1}G(\mathcal O_{X_S}/\mathcal O_{X_S}(-kD))
\]
的一个元素给出。于是可以定义局部 Hecke 叠
\[
\mathcal{H}ck_{G,\mathrm{Div}^d}=\big[L^+_{\mathrm{Div}^d}G\backslash\mathrm{Gr}_{G,\mathrm{Div}^d}\big]\longrightarrow\mathrm{Div}^d
\]
作为 $v$-叠。对满足 $|I|=d$ 的有限集 $I$，定义
\[
\mathcal{H}ck^I_G=\mathcal{H}ck_{G,\mathrm{Div}^d}\times_{\mathrm{Div}^d}(\mathrm{Div}^1)^I\longrightarrow(\mathrm{Div}^1)^I,
\]
其中 $(\mathrm{Div}^1)^I\to\mathrm{Div}^d$ 是求和映射。它在 $S\in\mathrm{Perf}_{\mathbb F_q}$ 上的取值，是由以下数据给出的群胚：
\begin{enumerate}[label=(\arabic*),leftmargin=2.4em,itemsep=0.2em]
\item $X_S$ 上的一组 1 次相对 Cartier 除子 $(D_i)_{i\in I}$，
\item $X_S$ 上的两个 $G$-丛 $\mathcal E$ 与 $\mathcal E'$，
\item 亚纯同构 $\mathcal E|_{X_S\setminus\bigcup_{i\in I}D_i}\xrightarrow{\ \sim\ }\mathcal E'|_{X_S\setminus\bigcup_{i\in I}D_i}$。
\end{enumerate}

\subsection{Satake 范畴（[9, Chapter VI.7.1]）}

设 $\Lambda$ 是挠的。Satake 范畴
\[
\mathrm{Sat}^I_G(\Lambda)
\]
定义为 $D_{\text{ét}}(\mathcal{H}ck^I_G,\Lambda)$ 中满足以下条件的复形 $A$ 的子范畴：
\begin{enumerate}[label=(\arabic*),leftmargin=2.4em,itemsep=0.2em]
\item $A$ 在 Hecke 叠上有界支撑，
\item $A$ 在 $(\mathrm{Div}^1)^I$ 上是 ULA 的，
\item $A$ 在 $(\mathrm{Div}^1)^I$ 上是平坦反常的（flat perverse）。
\end{enumerate}
\ysrc{148}
最后一个条件是说：对任意 $\Lambda$-模 $M$，$A\otimes^L_\Lambda M$ 是反常的。这里的反常性条件是说：
\begin{enumerate}[label=(\arabic*),leftmargin=2.4em,itemsep=0.2em]
\item 对任意由曲线上 $r$ 个互异点给出的态射 $x:\Spa(C,C^+)\to\mathcal{H}ck^I_G$，若它落在那 $r$ 个互异点处与 $\mu_1,\dots,\mu_r\in X_*(T)^+$ 相联系的 Schubert 胞腔中（这些 Schubert 胞腔照常由 $B^+_{dR}$-格的相对位置给出），则 $x^*A\in D^{\le-\sum_{i=1}^r\langle\mu_i,2\rho\rangle}$。
\item 对 $\mathbb D(A)$（其 Verdier 对偶）同样成立。
\end{enumerate}

最初的结果之一如下。令
\[
R\pi_{G*}:\mathrm{Sat}^I_G(\Lambda)\longrightarrow D_{\text{ét}}((\mathrm{Div}^1)^I,\Lambda)
\]
为拉回到 $\mathrm{Gr}_{G,\mathrm{Div}^d}\times_{\mathrm{Div}^d}(\mathrm{Div}^1)^I$ 再前推到 $(\mathrm{Div}^1)^I$ 的复合。
\begin{thm}{4.1}
函子 $R\pi_{G*}$ 取值于这样的复形 $C\in D_{\text{ét}}((\mathrm{Div}^1)^I,\Lambda)$：对一切 $i\in\mathbb Z$，$\mathcal H^i(C)$ 是有限投射 $\Lambda$-模构成的局部系统。函子
\[
\bigoplus_{i\in\mathbb Z}\mathcal H^i(R\pi_{G*}):\mathrm{Sat}^I_G(\Lambda)\longrightarrow\mathrm{LocSys}((\mathrm{Div}^1)^I,\Lambda)
\]
是正合、忠实且保守的。\label{thm:a3-4-1}
\end{thm}

有 $\mathrm{Div}^1=\mathrm{Spd}(\widehat E)/W_E$。这给出态射
\[
\mathrm{Div}^1\longrightarrow[*/W_E].
\]
可以证明下述结果，见 [9, Chapter IV.7]。\tn{底本 p.147 处记作 $\mathrm{Spd}(\breve E)$，此处作 $\mathrm{Spd}(\widehat E)$；二者指同一对象，照原样保留。}
\begin{thm}{4.2（Drinfeld 引理）}
对任意小 $v$-叠 $X$，拉回函子
\[
D_{\text{ét}}(X\times[*/W^I_E],\Lambda)\longrightarrow D_{\text{ét}}(X\times(\mathrm{Div}^1)^I,\Lambda)
\]
(1) 全忠实，且限制为局部系统上的等价；

(2) 当 $I$ 只有一个元素时是等价。\label{thm:a3-4-2}
\end{thm}

与通常的 shtuka 模空间的经典情形相比，这个证明出人意料地「简单」。这是我们的工作比有限域上函数域的「经典情形」更简单的少数场合之一。我们稍后会把该定理应用于 $X=\Bun_G$，但现在先记下下述推论。
\begin{cor}{4.3}
存在正合、忠实且保守的函子
\[
F^I:\mathrm{Sat}^I_G(\Lambda)\longrightarrow\Rep_\Lambda(W^I_E),
\]
其中 $\Rep_\Lambda(W^I_E)$ 是 $W_E$ 在有限型投射 $\Lambda$-模上的连续表示的范畴。\label{cor:a3-4-3}
\end{cor}

\subsection{卷积与融合}

\subsubsection{卷积（[9, Chapter VI.8]）}

有通常的卷积图表
\begin{center}
\begin{tikzcd}[column sep=large]
& {L^+G\backslash LG\times^{L^+G}LG/L^+G} \arrow[dl, "a"'] \arrow[dr, "b"] & \\
{\mathcal{H}ck_{G,\mathrm{Div}^d}\times_{\mathrm{Div}^d}\mathcal{H}ck_{G,\mathrm{Div}^d}} & & {\mathcal{H}ck_{G,\mathrm{Div}^d}}
\end{tikzcd}
\end{center}
